\documentclass[12pt,oneside,a4paper]{book}

%--------------------------------------
% Encodage et langue
%--------------------------------------
\usepackage[utf8]{inputenc}
\usepackage[T1]{fontenc}
\usepackage[french]{babel}
\usepackage[autolanguage]{numprint}

%--------------------------------------
% Règles de césure
%--------------------------------------
\usepackage{hyphenat}
\hyphenation{matematica recuperar}

%--------------------------------------
% Graphiques
%--------------------------------------
\usepackage{graphicx}
\usepackage[table]{xcolor}
\graphicspath{{./pictures/}{./images/}}

%--------------------------------------
% Mise en page
%--------------------------------------
\usepackage[left=25mm,right=25mm,top=25mm,bottom=25mm,paper=a4paper]{geometry}
\usepackage{setspace}
\onehalfspacing

%--------------------------------------
% Tables des matières par chapitre
%--------------------------------------
%--------------------------------------
% Divers
%--------------------------------------
\usepackage{fancyhdr}
\usepackage{multicol}
\usepackage{multirow}
\usepackage{amsfonts}
\usepackage{amsmath}
\usepackage{caption}
\usepackage{longtable}
\usepackage{array}
\usepackage{enumitem}
\usepackage{amssymb}
\usepackage{pdfpages}
\usepackage[Glenn]{fncychap}
\usepackage{etoolbox}
\usepackage{ragged2e}
\usepackage{float}
\usepackage{url}
\usepackage{silence}
\usepackage{mdframed}

% Packages requis pour les graphiques de convergence IA
\usepackage{tikz}
\usepackage{pgfplots}
\pgfplotsset{compat=1.18} % Assure la compatibilité des axes

%--------------------------------------
% Hyperliens
%--------------------------------------
\usepackage[
    colorlinks=true,
    linkcolor=black,
    urlcolor=blue,
    citecolor=black
]{hyperref}
\usepackage[french]{minitoc}
\makeatletter
\renewcommand{\mtcPackageWarningNoLine}[3][]{\ifnum\pdfstrcmp{#2}{minitoc(hints)}=0\relax\else\PackageWarningNoLine{#2}{#1\MessageBreak #3}\fi}
\makeatother

\frenchbsetup{StandardLists=true}

%%%%%%%%%%%%%%%%%%%%%%%%%%%%%%%%%%
\newcommand\myemptypage{
    \null
    \thispagestyle{empty}
    \addtocounter{page}{-1}
    \newpage
}
%%%%%%%%%%%%%%%%%%%%%%%%%%%%%%%%%%
\newcommand{\ac}[1]{#1}
\setlength{\headheight}{28pt}
\setlength{\emergencystretch}{2em}
\hbadness=10000
\vbadness=10000
\hfuzz=20pt
\vfuzz=70pt

\begin{document}

\dominitoc
\setcounter{tocdepth}{3}
\setcounter{secnumdepth}{3}

\renewcommand\thepage{\arabic{page}}
\setcounter{page}{3}
\thispagestyle{empty}

% Page de garde PDF
\includepdf[pages=1]{files/page_garde.pdf}

\newpage

% Pages préliminaires
\include{front/Garde}
\include{front/Remerciement}
\include{front/Dedicace1}
\include{front/Dedicace2}
\include{front/Acronymes}

\makeatletter
\renewcommand{\l@chapter}[2]{%
  \addpenalty{-\@highpenalty}%
  \vskip 1.0em \@plus\p@
  \setlength\@tempdima{1.5em}%
  \begingroup
    \parindent \z@
    \rightskip \@tocrmarg
    \parfillskip -\rightskip
    \leavevmode
    \advance\leftskip\@tempdima
    {\bfseries Chapitre #1}
  \endgroup
}
\makeatother

\tableofcontents
\hfill

\listoffigures
\addcontentsline{toc}{section}{\hspace*{-1.5em}\textbf{Table des figures}}
\hfill

\listoftables
\addcontentsline{toc}{section}{\hspace*{-1.5em}\textbf{Liste des tableaux}}
\hfill

\renewcommand{\contentsname}{Table des mati\`eres}
\thispagestyle{empty}
\pagestyle{fancy}

\fancyfoot[C]{\thepage}
\fancyhead[R]{\rightmark}
\fancyhead[L]{}
\input{restructuration/notice}
%%%%% INTRO GENERALE %%%%%
\chapter*{Introduction Générale}
\addcontentsline{toc}{section}{\hspace*{-1.5em}\textbf{Introduction Générale}}
\markboth{Introduction Générale}{Introduction Générale}
\input{front/IntGenerale}

\renewcommand{\arraystretch}{1.3}

% Chapitres
\include{chapters/Chapter1}
\include{chapters/Chapter2}
\input{restructuration/methodologie}
% \include{chapters/Chapter3}
% \include{chapters/Chapter4}
% \include{chapters/Chapter5}

\input{restructuration/release1}
\input{restructuration/release2}
\include{comparaison}

% \addcontentsline{toc}{section}{\hspace*{-1.5em}\textbf{Conclusion Générale}}
% \renewcommand\headrulewidth{1pt}
% \fancyhead[R]{Conclusion Générale}
% \include{front/ConclusionGenerale}

% Bibliographie
\renewcommand{\bibname}{Webographie}
\cleardoublepage
\phantomsection
% \addcontentsline{toc}{section}{\hspace*{-1.5em}\textbf{Webographie}}
\markboth{Webographie}{Webographie}
\fancyhead[R]{\rightmark}

\bibliographystyle{unsrt}
\bibliography{biblio}

\newpage
\addcontentsline{toc}{chapter}{\large{Abstract}}
\renewcommand\headrulewidth{1pt}
\fancyhead[R]{Abstract}

% IMPORTANT :
% Décommente cette ligne seulement si le fichier Abstract.tex existe
\input{front/Abstract}

% Dernière page PDF
% \includepdf[pages=1]{files/plagiatPage.pdf}

\end{document}